\documentclass[11pt]{article}
\usepackage[letterpaper, margin=2cm]{geometry}
\usepackage{titlesec}
\usepackage{mdframed}
\usepackage[dvipsnames]{xcolor} % for color names, must be loaded before tikz
\usepackage{ifthen}
\usepackage{fancyhdr}
\usepackage{comment}
\usepackage{titling}
\usepackage{hyperref}
\usepackage{enumitem}
\usepackage{tikz}
\usepackage{amsmath, amssymb, amsthm}
\usepackage{mathtools}
\usepackage[braket, qm]{qcircuit}

\newcommand{\abs}[1]{\left\lvert #1\right\rvert}


\providecommand{\due}{Due Wednesday, September 11th at 11:59 PM}
\lhead{CS358H, M375T, ES377, PHY341} \rhead{}
\lfoot{\due} \cfoot{} \rfoot{Page \thepage}
\renewcommand{\footrulewidth}{0.4pt}
\pagestyle{fancy}

% Eliminates the spacing in the title that remains from the empty author section.
\preauthor{}
\postauthor{}

\titleformat{\section}[runin]{\large\bfseries}{\thesection .}{3pt}{}
\titleformat{\subsection}[runin]{\bfseries}{\thesubsection)}{3pt}{}
\renewcommand\thesubsection{\alph{subsection}}

% Defines the solution environment. Toggle solutions between true and false to either show or hide solutions. Also, the solution environment takes an optional argument of arbitrary text to be inserted in the solution header.
\newboolean{solutions}
\setboolean{solutions}{true}
\ifthenelse{\boolean{solutions}}
{\newenvironment{solution}{\begin{mdframed}[skipbelow=0pt, linecolor=White, backgroundcolor=Green!10]\textbf{Solution:}}{\end{mdframed}}}
{\excludecomment{solution}}

\DeclareMathOperator{\CNOT}{CNOT}
\DeclareMathOperator{\Tr}{Tr}
\DeclareMathOperator{\Ev}{\mathbb{E}}

\allowdisplaybreaks

\begin{document}

\title{Introduction to Quantum Information Science\\Homework 1}
\date{\due}

\maketitle

\section{Stochastic and Unitary Matrices.} 

\subsection{[8 Points]} Of the following matrices, which ones are stochastic? Which
ones are unitary?

\begin{gather*}
A =
\begin{bmatrix}
1 & 0\\
0 & 0
\end{bmatrix}\!,\;
B=
\begin{bmatrix}
0 & 1\\
1 & 0
\end{bmatrix}\!,\;
C=
\begin{bmatrix}
1 & \frac{2}{3}\\[0.1em]
0 & \frac{1}{3}
\end{bmatrix}\!,\;
D=
\begin{bmatrix}
1 & 0\\
0 & i
\end{bmatrix}\!,\\
E =
\begin{bmatrix}
2 & \frac{1}{2}\\[0.1em]
-1 & \frac{1}{2}
\end{bmatrix}\!,\;
F = \frac{1}{\sqrt{2}}
\begin{bmatrix}
1 & 1\\
1 & -1
\end{bmatrix}\!,\;
G=
\begin{bmatrix}
\frac{3}{5} & \frac{4}{5}\\[0.1em]
\frac{4}{5} & -\frac{3}{5}
\end{bmatrix}\!,\;
H=
\begin{bmatrix}
\frac{3i}{5} & \frac{4}{5}\\[0.1em]
\frac{4}{5} & -\frac{3i}{5}
\end{bmatrix}
\end{gather*}

\begin{solution}
	\begin{enumerate}
		\item A: A is neither stochastic, nor unitary. 
		\item B: B is stochastic. $\overline{B}^T B = \begin{bmatrix}
			0 & 1 \\ 1 & 0
		\end{bmatrix}\begin{bmatrix}
			0 & 1 \\ 1 & 0
		\end{bmatrix} = B\overline{B}^T = I$. B is unitary. 
		\item C: C is stochastic. $\overline{C}^T C = \begin{bmatrix}
			1 & 0 \\  \frac{2}{3} & \frac{1}{3}
		\end{bmatrix}\begin{bmatrix} 
			1 & \frac{2}{3} \\ 0 & \frac{1}{3}
		\end{bmatrix} = \begin{bmatrix}
			1 & \frac{2}{3} \\ \frac{2}{3} & \frac{5}{9}
		\end{bmatrix} \neq I$. Thus $C$ is not a unitary matrix.   
		\item D: D is not stochastic. $\overline{D}^TD = \begin{bmatrix}
			1 & 0 \\ 0 & -i
		\end{bmatrix}\begin{bmatrix}
			1 & 0 \\ 0 & i
		\end{bmatrix} = \begin{bmatrix}
			1 & 0 \\ 0 & 1
		\end{bmatrix} = \begin{bmatrix}
			1 & 0 \\ 0 & i
		\end{bmatrix}\begin{bmatrix}
			1 & 0 \\ 0 & -i
		\end{bmatrix} = D \overline{D}^T$. Thus $D$ is unitary.     
		\item E: E is not stochastic. $\overline{E}^T E = \begin{bmatrix}
			2 & -1 \\ \frac{1}{2} & \frac{1}{2}
		\end{bmatrix}\begin{bmatrix}
			2 & \frac{1}{2} \\ -1 & \frac{1}{2}
		\end{bmatrix} = \begin{bmatrix}
			5 & \frac{1}{2} \\  \frac{1}{2} & \frac{1}{2}
		\end{bmatrix} \neq I$. E is not unitary 
		\item F: F is not stochastic. $\overline{F}^T F = \frac{1}{2} \begin{bmatrix}
			1 & 1 \\ 1 & -1
		\end{bmatrix}\begin{bmatrix}
			1 & 1 \\ 1 & -1
		\end{bmatrix} = \frac{1}{2}\begin{bmatrix}
			2 & 0 \\ 0 & 2
		\end{bmatrix} = I = \frac{1}{2} \begin{bmatrix}
			1 & 1 \\ 1 & -1
		\end{bmatrix}\begin{bmatrix}
			1 & 1 \\ 1 & -1
		\end{bmatrix} = F \overline{F}^T$. F is unitary.
		\item G: G is not stochastic. G is symmetric so $G^T = G$ and since G is real $\overline{G} = G$. $G \overline{G}^T = \overline{G}^T G = \frac{1}{25} \begin{bmatrix}
			3 & 4 \\ 4 & -3
		\end{bmatrix}\begin{bmatrix}
			3 & 4 \\ 4 & -3
		\end{bmatrix} = \frac{1}{25}\begin{bmatrix}
			25 & 0 \\ 0 & 25
		\end{bmatrix} = I$. G is unitary.  
		\item H: H is not stochastic. H is symmetric so $H^T = H$. $\overline{H}^T H = \overline{H} H = \frac{1}{25}  \begin{bmatrix}
			-3i & 4 \\ 4 & 3i
		\end{bmatrix}\begin{bmatrix}
			3i & 4 \\ 4 & -3i
		\end{bmatrix} = \frac{1}{25}\begin{bmatrix}
			25 & -24i \\ 24i & 25
		\end{bmatrix} \neq I$. H is not unitary.  
	\end{enumerate}
\end{solution}


\subsection{[3 Points]} Show that any stochastic matrix that is also unitary must be a permutation matrix. 

\begin{solution}
	Suppose $A \in M_{n \times n}$ is a stochastic, unitary matrix. Since $A$ is stochastic it is a real, nonnegative matrix. Thus A is orthogonal, since the conjugate of real matrices is the matrix itself. Let $a_{i, j}$ be the entry on the $i$th row and $j$th column. Note $\forall j \in [1, \dots, n]$, $\sum_{i = 1}^{n} a_{i, j} = 1$ and $a_{i, j} \geq 0$ for all $i, j \in [1, \dots, n]$. 
	
	Looking at the diagonal entris of $A^T A = I$, we see.  
	\[\sum_{i = 1}^n a^2_{i, j} = 1\]

	\[1 = \left(\sum_{i = 1}^{n} a_{i, j}\right)^2 = \sum_{i = 1}^n a^2_{i, j} + 2\sum_{r = 1}^{n - 1} \sum_{s = r + 1}^{n} a_{r, j} a_{s, j} = 1 + 2\sum_{r = 1}^{n - 1} \sum_{s = r + 1}^{n} a_{r, j} a_{s, j} \implies \sum_{r = 1}^{n - 1} \sum_{s = r + 1}^{n} a_{r, j} a_{s, j} = 0\]'

	Since $a_{i, j} \geq 0 \implies $ either $a_{r, i} = 0$ or $a_{s, i} = 0$ for all $r \in [1, n - 1]$ and $s \in [r, n]$. So we have three cases: 

	\begin{enumerate}
		\item Every entry in the column is zero. This is a contradiction since $A$ is orthogonal and thus invertible. Thus if every entry in the column is zero then A cannot be invertible. By contradiction, this case is not true. 
		\item Only one entry in the column is 1, everything else is zero. Thus the sum of the elements on the column is 1 and the product of any pair of entries is 0.   
		\item More than one entry is nonzero: Let $a_{r, i}$ and $a_{s, i}$ be two such entries. Then $a_{r, i} a_{s, i} \neq 0$ which is a contradiction. Thus by contradiction, this cases is impossible.  
	\end{enumerate}

	Thus only one entry in the column is a 1. 

	So our matrix is a binary matrix and we have $n$ ones. We want to show every row has exactly 1 one. 
	
	Assume the contrary $\exists$ a row $k$ row with no one. This is equivalent to saying there is a row with multiple ones since there are $n$ rows and $n$ ones if there is a row with no one then there must be a row with multiple ones. Then all entries on the row must be 0 which causes the matrix to not be invertible. This is a contradiction since $A \in U(n)$ thus $A \in GL_n(\mathbb{C})$. Thus by contradiction every row must have exactly 1 one.
	
	Thus the matrix A is a binary matrix where every row and column has exactly 1 one. Thus by definition of permutation matrices, A is a permutation matrix.  
	
\end{solution}


\subsection{[1 Point]} Stochastic matrices preserve the $1$-norms of nonnegative vectors, while
unitary matrices preserve $2$-norms. \ Give an example of a $2\times
2$\ matrix, other than the identity matrix, that preserves the $4$-norm of
real vectors $\begin{bsmallmatrix} a\\b\end{bsmallmatrix}$: that is, $a^{4}+b^{4}$.

\begin{solution}
	An example of $2 \times 2$ that preseves 4 norm is $\begin{bmatrix}
		0 & -1 \\ 1 & 0
	\end{bmatrix}$. 
	
	$\forall a, b \in \mathbb{R}$, $\begin{bmatrix}
		0 & -1 \\ 1 & 0
	\end{bmatrix}\begin{bmatrix}
		a \\ b
	\end{bmatrix} = \begin{bmatrix}
		-b \\ a
	\end{bmatrix}$. $\abs{\abs{\begin{bmatrix}
		-b \\ a
	\end{bmatrix}}}_{4} = b^4 + a^4 = a^4 + b^4 = \abs{\abs{\begin{bmatrix}
		a \\ b
	\end{bmatrix}}}_4$   
\end{solution}

\subsection{[Extra credit, 4 Points]} Give a characterization of all real matrices that preserve the
$4$\textit{-norms} of real vectors. \ Hopefully, your characterization will
help explain why preserving the $2$-norm, as quantum mechanics does, leads to
a much richer set of transformations than preserving the $4$-norm does.



\section{Tensor Products}

\subsection{[1 Point]} Calculate the tensor product 
\[
\begin{bmatrix}
\frac{2}{3}\\[0.1em]
\frac{1}{3}
\end{bmatrix}\otimes
\begin{bmatrix}
\frac{1}{5}\\[0.1em]
\frac{4}{5}
\end{bmatrix}.
\]

\begin{solution}
	\[\begin{bmatrix}
		\frac{2}{3}\\
		\frac{1}{3}
		\end{bmatrix}\otimes
		\begin{bmatrix}
		\frac{1}{5}\\
		\frac{4}{5}
		\end{bmatrix} = \begin{bmatrix}
			\frac{2}{15} \\[0.1em]
			\frac{8}{15} \\[0.1em]
			\frac{1}{15} \\[0.1em]
			\frac{4}{15} 
		\end{bmatrix}\]
\end{solution}



\subsection{[5 Points]} Of the following length-$4$ vectors, decide which ones are factorizable as a tensor product of two $2 \times 1$ vectors, and factorize them. \ (Here the vector entries
should be thought of as labeled by $00$, $01$, $10$, and $11$ respectively.)%
\[
A=
\begin{bmatrix}
	\frac{2}{9}\\[0.1em]
	\frac{1}{9}\\[0.1em]
	\frac{4}{9}\\[0.1em]
	\frac{2}{9}
\end{bmatrix}\!,\;
B=
\begin{bmatrix}
	0\\
	1\\
	0\\
	0
\end{bmatrix}\!,\;
C=
\begin{bmatrix}
	\frac{1}{4}\\[0.1em]
	\frac{1}{4}\\[0.1em]
	\frac{1}{4}\\[0.1em]
	\frac{1}{4}
\end{bmatrix}\!,\;
D=
\begin{bmatrix}
	0\\
	\frac{1}{2}\\[0.1em]
	\frac{1}{2}\\
	0
\end{bmatrix}\!,\;
E=
\begin{bmatrix}
	0\\
	\frac{1}{2}\\
	0\\
	\frac{1}{2}
\end{bmatrix}.
\]

\begin{solution}
	\begin{enumerate}
		\item A: \[
			\begin{bmatrix}
				\frac{2}{9}\\[0.1em]
				\frac{1}{9}\\[0.1em]
				\frac{4}{9}\\[0.1em]
				\frac{2}{9}
			\end{bmatrix} = \begin{bmatrix}
				\frac{1}{3} \\[0.1em] 
				\frac{2}{3}
			\end{bmatrix} \otimes 
			\begin{bmatrix}
				\frac{2}{3} \\[0.1em]
				\frac{1}{3}
			\end{bmatrix}
		\]
		\item B:
		
		\[
			\begin{bmatrix}
				0 \\ 1 \\ 0 \\ 0
			\end{bmatrix} = \begin{bmatrix}
				1 \\ 0
			\end{bmatrix} \otimes 
			\begin{bmatrix}
				0 \\ 1
			\end{bmatrix}
		\]
		
		\item C: \[
			\begin{bmatrix}
				\frac{1}{4}\\[0.1em]
				\frac{1}{4}\\[0.1em]
				\frac{1}{4}\\[0.1em]
				\frac{1}{4}
			\end{bmatrix} = \begin{bmatrix}
				\frac{1}{2} \\[0.1em] 
				\frac{1}{2}
			\end{bmatrix} \otimes 
			\begin{bmatrix}
				\frac{1}{2} \\[0.1em]
				\frac{1}{2}
			\end{bmatrix}
		\]
		\item D: D cannot be factored. 
		\item E: \[
			\begin{bmatrix}
				0\\[0.1em]
				\frac{1}{2}\\[0.1em]
				0\\[0.1em]
				\frac{1}{2}
			\end{bmatrix} = \begin{bmatrix}
				\frac{1}{2} \\[0.1em] 
				\frac{1}{2}
			\end{bmatrix} \otimes 
			\begin{bmatrix}
				0 \\
				1
			\end{bmatrix}
		\]
	\end{enumerate}
	
\end{solution}


\subsection{[3 Points]} Prove that there's no $2 \times 2$ \emph{real} matrix $A$ such
that
\[
A^2=
\begin{bmatrix}
	1 & 0\\
	0 & -1
\end{bmatrix}\!.
\]
This observation perhaps helps to explain why the complex numbers play such a
central role in quantum mechanics.

\begin{solution}
	Assume the contrary that there exists $A = \begin{bmatrix}
		a & b \\ c & d
	\end{bmatrix} \in M_{2 x 2}(\mathbb{R})$ such that $A^2 = \begin{bmatrix}
		1 & 0 \\ 0 & -1
	\end{bmatrix}$. 

	Then, \[
		\begin{bmatrix}
			a & b \\ c & d
		\end{bmatrix}  \begin{bmatrix}
			a & b \\ c & d
		\end{bmatrix} = \begin{bmatrix}
			a^2 + bc & ab + bd \\ ac + dc & bc + d^2 
		\end{bmatrix} = \begin{bmatrix}
			1 & 0\\
			0 & -1
		\end{bmatrix}
	\]

	Thus we have 4 equations: 

	\[a^2 + bc = 1\]
	\[ab + bd = b(a + d) = 0\]
	\[ac + dc = c(a + d) = 0\]
	\[bc + d^2 = -1\]

	$bc = 1 - a^2 \implies 1 - a^2 + d^2 = 1 \implies 2 = a^2 - d^2 = (a + d)(a - d)$. 
	
	$b(a + d) - 0 \implies a + d = 0$ or $b = 0$. Assume the contrary that $a + d = 0$, then $(a + d)(a - d) = a^2 - d^2 = 0 \neq 2$ which is a contradiction. Thus $b = 0$. Thus $bc + d^2 = d^2 = -1$. $\forall r \in \mathbb{R}$, $r^2 \geq 0$. Thus $\not \exists d \in \mathbb{R} \ni d^2 = -1 < 0$. Thus $d \not \in \mathbb{R}$. Thus $A \not \in M_{2 \times 2}(\mathbb{R})$. Thus a contradiction occurs. Thus there exists no 2x2 real matrix whose square is $ \begin{bmatrix}
		1 & 0 \\ 0 & -1
	\end{bmatrix}$.                 
\end{solution}

\section{Dirac Notation}
 
\subsection{[2 Points]} Let $\ket{\psi} = \frac{\ket{0} + 2 \ket{1}}{\sqrt{5}}$ and $\ket{\phi} = \frac{2i\ket{0} + 3\ket{1}}{\sqrt{13}}$. What's $\ip{\psi}{\phi}$?

\begin{solution}
	\begin{align*}
		\ip{\psi}{\phi} &= \left( \frac{\overline{1}\bra{0} + \overline{2}\bra {1}}{\sqrt{5}} \right) \left(\frac{2i\ket{0} + 3\ket{1}}{\sqrt{13}} \right) \\
		&= \frac{1}{\sqrt{65}}(2i + 6) = \frac{1}{\sqrt{65}}(6 + 2i)
	\end{align*}
\end{solution}


\subsection{[1 Point]} Usually quantum states are normalized: $\ip{\psi}{\psi} = 1$. The state $\ket{\varphi} = 2i \ket{0} - 3i \ket{1}$ is not normalized. What constant $A$ makes $\ket{\psi} = \frac{\ket{\varphi}}{A}$ a normalized state?

\begin{solution}
	\begin{align*}
		A &= \sqrt{\ip{\varphi}{\varphi}} \\
		&= \sqrt{(-2i \bra{0} + 3i\bra{1}) (2i \ket{0} - 3i\ket{1})} \\
		&= \sqrt{4 + 9} = \sqrt{13}
	\end{align*} 

	\begin{align*}
		\ip{\frac{\varphi}{A}}{\frac{\varphi}{A}} &= \frac{1}{A^2} \ip{\varphi}{\varphi} \\ 
		&= \frac{1}{13} * 13 = 1
	\end{align*}
\end{solution}


\subsection{[2 Points]} Define $\ket{i} = \frac{\ket{0} + i \ket{1}}{\sqrt{2}}$ and $\ket{-i} = \frac{\ket{0} - i\ket{1}}{\sqrt{2}}$. Show (explicitly or implicitly) that the vectors $\ket{i}$ and $\ket{-i}$ form an orthonormal basis for $\mathbb{C}^2$.
(Hint: show that any vector in $\mathbb{C}^2$ can be decomposed as a linear combination of $\ket{i}$ and $\ket{-i}$.)

\begin{solution}
	A basis of a vector space is a collection of linear independent vectors which spans the vectors. 

	\begin{enumerate}
		\item $\ket{i}$ and $\ket{-i}$ are linearly independent: Suppose $\exists a, b \in \mathbb{C} \ni a \ket{i} + b \ket{-i} = 0 \implies a \left(\frac{\ket{0} + i\ket{1}}{\sqrt{2}}\right) + b \left(\frac{\ket{0} - i\ket{1}}{\sqrt{2}}\right) = \frac{1}{\sqrt{2}} \left((a + b)\ket{0} + (a - b) i \ket{1}\right) = 0 \implies a + b =0$ and $(a - b) i = 0 \implies a - b = 0$. Thus $a = -b$ and $a = b$. Thus $a = -a \implies a = 0 \implies b = 0$. Thus $\ket{i}$ and $\ket{-i}$. Are linearly independent. 
		\item If a collection of vectors spans a vectorspace iff if the dimension of the span is equal to the dimension of the vectorspace. The dimension of the span of a collection of linearly independent vectors in equal to the number of vectors in the collection. Thus $\dim \text{span}(\ket{i}, \ket{-i}) = 2 = \dim \mathbb{C}^2$. Thus $\left\{ \ket{i}, \ket{-i} \right\}$ spans $\mathbb{C}^2$. 
	\end{enumerate}

	Thus $\left\{ \ket{i}, \ket{-i} \right\}$ is a basis of $\mathbb{C}^2$. 
	
	A basis is orthonormal iff if all pairs of vectors in the basis collection have a inner product equal to 0. 

	\begin{align*}
		\ip{i}{-i} &= \frac{1}{2}(\bra{0} - i \bra{1})(\ket{0} - i \ket{1})\\
		&= \frac{1}{2}(1 + i^2) = 0
	\end{align*}

	\begin{align*}
		\ip{-i}{i} = \overline{\ip{i}{-i}} = 0
	\end{align*}

	Thus $\left\{ \ket{i}, \ket{-i} \right\}$ is orthonormal  

	Thus $\left\{ \ket{i}, \ket{-i} \right\}$ is a orthonormal basis of $\mathbb{C}^2$ 
	
\end{solution}


\subsection{[2 Points]} Write the normalized vector $\ket{\psi}$ from part (b) in the $\{\ket{i}, \ket{-i}\}$-basis.

\begin{solution}
	$A = \frac{1}{\sqrt{2}} \begin{bmatrix}
		1 & 1 \\ i & -i 
	\end{bmatrix}$ is the matrix which turns a vector in the $\left\{ \ket{i}, \ket{-i} \right\}$ basis to the standard basis $\left\{ \ket{0}, \ket{1} \right\}$. So $A^{-1} = \frac{1}{\sqrt{2}} \begin{bmatrix}
		1 & -i \\ 1 & i
	\end{bmatrix}$ is the matrix that takes a vector in the $\left\{ \ket{0}, \ket{1} \right\}$ basis to its equivalent in $\left\{ \ket{i}, \ket{-i} \right\}$ basis. 
	
	So $\ket{\psi} = \frac{1}{\sqrt{26} } \left((2i - 3) \ket{i} + (2i + 3) \ket{-i}\right)$. 
	

	Verification: 
	\begin{align*}
		\frac{1}{\sqrt{26} } \left((2i - 3) \ket{i} + (2i + 3) \ket{-i}\right) &= \frac{1}{2\sqrt{13}}\left((2i - 3)(\ket{0} + i\ket{1}) + (2i + 3)(\ket{0} - i\ket{1})\right) \\
		&= \frac{1}{2\sqrt{13}}(4i \ket{0} - 2 \ket{1} + 2 \ket{1} - 3\ket{0} + 3\ket{0} - 6i \ket{1}) \\
		&= \frac{1}{\sqrt{13}}(2i \ket{0} - 3i\ket{1}) = \ket{\psi}
	\end{align*}
\end{solution}

\end{document}

